\chapter{Comparison Interfaces: \texttt{std::map} and \texttt{std::vector}}\label{app:interfaces}

% [AP-G5-Anhang / AP-EN] Mirror of anhang/de/F_comparison_interfaces.tex. Full standard interface tables (C++23).
This appendix presents --- complementing Section~\ref{sec:measurement-system} --- the complete
listing of all standard interface functions of the two comparison hulls \texttt{std::map} (ordered
associative) and \texttt{std::vector} (sequence) as of C++23, with their officially expected
semantics and their complexity guarantee per the ISO C++ standard~\cite{iso_cpp}. It also serves as
the central reference for the further development of the cache engine: because the search-algorithm
hull exchanges its internals at compile time --- each axis choice yields a compiled binary of its own
---, while at run time only the parameters of the dynamic sub-axes are readjusted (the registered
dynamic dimensions of the measurement realm, Section~\ref{sec:measurement-system}), it becomes
necessary to explain, per standard function, which behaviour the outer contract guarantees.

\emph{Variadic mapping.} The hybrid search engine of the candidate under test PRT-ART ties the number of
parameters to the hull: one template parameter demands all function interfaces of a
\texttt{std::vector}, two those of a \texttt{std::map}, and $N>2$ those of a
\texttt{std::map<K,\,std::tuple<\ldots>>} (Section~\ref{sec:impl-architecture}). The cache engine
derives the types from the same rule --- with one parameter the key is assigned implicitly as a
monotonically counted 64-bit value, with $N>2$ the value becomes the tuple ---; its two contract
hulls are implemented as separate family adapters: the SearchAlgorithm family (concept Map) via the
functional drive of the
module boundary with the five operations \texttt{tier\_insert}, \texttt{tier\_lookup},
\texttt{tier\_erase}, \texttt{tier\_clear}, and \texttt{tier\_size}, the Sequence family via the
sequence drive with \texttt{tier\_push\_back}, \texttt{tier\_at}, \texttt{tier\_size}, and
\texttt{tier\_clear}; an adapter cutting across the concept is ruled out by the type system. The case
$N>2$ is type derivation without an adapter of its own; the module boundary carries key and value
throughout as 64-bit values. The shared functions (iteration,
\texttt{size}/\texttt{empty}/\texttt{clear}, \texttt{swap}, \texttt{get\_allocator}) have the same
name and the same role in both hulls; of these, both hulls pass \texttt{size} and \texttt{clear}
across the module boundary, \texttt{empty} is derived host-side, iteration and \texttt{swap} do not
cross the boundary, and \texttt{get\_allocator} is, on the map side, an optional proxy sub-interface
outside the oracle (Section~\ref{asec:if-decomp}).

\emph{Composite functions.} Some functions arise internally as a composition of more primitive
operations --- for instance \texttt{operator[]} as \texttt{lower\_bound} followed by default
insertion on a missing key, or \texttt{at} as \texttt{find} with a thrown
\texttt{std::out\_of\_range} exception on absence ---; such cases are marked in the
\enquote{Note} column. The \enquote{Since} column names the introducing standard version. Three
terms of the tables stem from the more recent development of the standard: the node
handle\footnote{Since C++17, a node of an associative container can be unlinked with
\texttt{extract} as a node handle, without copy or move, and spliced into another container with
\texttt{insert} or \texttt{merge}~\cite{talbot2016splicing}.}, the deduction guides
(CTAD)\footnote{\emph{Class template argument deduction}, since C++17: the template arguments of a
class template are deduced from the constructor arguments, via guides formed implicitly from the
constructors or stated explicitly~\cite{spertus2016ctad}.}, and the transparent key
(Section~\ref{asec:if-map}). The exact cache-engine-internal function hops
(\emph{function-handle hops}\footnote{Project-internal term of the cache engine for the sequence of
internal function hops (handles) from the entry at the module boundary through the axis organs to
the result; the developer document records them per standard function and marks functions that
are not wired up as such.}) are documented in the developer document
\texttt{docs/\allowbreak{}ENTWICKLER-\allowbreak{}FUNCTION-\allowbreak{}HANDLE-\allowbreak{}HOPS.md}
of the cache engine; it records, per standard function, the hop path and the status at the module
boundary.

\emph{Conformance subset.} Before each measurement, the conformance gate
(Section~\ref{sec:measurement-system}) checks, per family, a core subset of the hull operations
against an oracle of the standard library --- \texttt{std::map} for the SearchAlgorithm family
(concept Map), \texttt{std::deque} for the Sequence family ---, so that only comparable things are
compared; the subset is marked with $\dagger$ in both tables. Under \texttt{--debug} the gate
additionally runs the standard test set of the tier interface (family-generic part and genus-specific
part).

\section{\texttt{std::map} (Ordered Associative)}\label{asec:if-map}
Table~\ref{tab:if-map} lists the complete interface of
\texttt{std::map<Key,\,T,\,Compare,\,Allocator>} as of C++23. A $\dagger$ marks the core operations
that the conformance gate (Section~\ref{sec:measurement-system}) checks before each measurement,
so that only comparable things are compared. Complexities refer to $n=\texttt{size()}$; for range
operations $N$ is the number of input elements; version markers and complexities follow the
standard text~\cite{iso_cpp} and are cross-checked against the online
reference~\cite{cppreference_map}. The seventeen marked operations are exactly the operation set of
the gate's twelve core check blocks: no marked operation remains unchecked, no checked operation is
unmarked (\texttt{size} is checked as well within the blocks, but belongs to the randomized core
block and is therefore not marked). Four of the marked operations correspond directly to operations
of the module boundary (\texttt{insert\_or\_assign} to the insert drive, \texttt{find} to the exact
lookup, \texttt{erase} and \texttt{clear} to their drives); \texttt{at}, \texttt{operator[]},
\texttt{count}, \texttt{contains}, \texttt{empty} as well as \texttt{insert}, \texttt{emplace}, and
\texttt{try\_emplace} are composed host-side --- \texttt{empty} from the size, the others from
exact lookup and insertion; the three non-overwriting forms run over one path and are set against
three different oracle calls; \texttt{begin}/\texttt{end}, \texttt{lower\_bound},
\texttt{upper\_bound}, and \texttt{equal\_range} are synthesized by the host in order via the
additive scan capability of the module boundary, and if a loaded binary lacks it, the gate logs
these blocks as skipped, not as passed; \texttt{key\_comp} is established at the oracle (key order
less-than on 64-bit values), because the hull carries no comparator across the boundary.

\begin{longtable}{@{}>{\raggedright\arraybackslash}p{3.1cm} >{\raggedright\arraybackslash}p{5.7cm} >{\raggedright\arraybackslash}p{2.7cm} >{\raggedright\arraybackslash}p{1.1cm}@{}}
  \caption{Complete \texttt{std::map} interface (C++23)}\label{tab:if-map}\\
  \toprule
  \textbf{Function} & \textbf{Official semantics} & \textbf{Complexity} & \textbf{Since} \\
  \midrule
\endfirsthead%
  \toprule
  \textbf{Function} & \textbf{Official semantics} & \textbf{Complexity} & \textbf{Since} \\
  \midrule
\endhead%
  \multicolumn{4}{@{}l}{\textbf{Construction, assignment, allocator}}\\
  \addlinespace[2pt]
  \texttt{(constructor)} & empty; from \texttt{Compare}/\texttt{Allocator}; range; copy/move (incl.\ allocator-extended); \texttt{initializer\_list}; \texttt{from\_range} since C++23 & $O(1)$ to $O(N\log N)$ & C++98 \\
  \texttt{(destructor)} & destroys all elements and nodes & $O(n)$ & C++98 \\
  \texttt{operator=} & copy/move/\texttt{initializer\_list} assignment with allocator propagation & $O(n)$ & C++98 \\
  \texttt{get\_allocator} & copy of the associated allocator & $O(1)$ & C++98 \\
  \addlinespace[3pt]
  \multicolumn{4}{@{}l}{\textbf{Element access}}\\
  \addlinespace[2pt]
  \texttt{at}\,$\dagger$ & reference to the mapped value; throws \texttt{std::out\_of\_range} on absence & $O(\log n)$ & C++11 \\
  \texttt{operator[]}\,$\dagger$ & reference; inserts a value-initialized element on absence (no \texttt{std::out\_of\_range}; the insertion can throw \texttt{std::bad\_alloc} or constructor exceptions) & $O(\log n)$ & C++98 \\
  \addlinespace[3pt]
  \multicolumn{4}{@{}l}{\textbf{Iterators}}\\
  \addlinespace[2pt]
  \texttt{begin}/\texttt{end}\,$\dagger$ & iterator to first element / end, in key order (\texttt{cbegin}/\texttt{cend} since C++11) & $O(1)$ & C++98 \\
  \texttt{rbegin}/\texttt{rend} & reverse iterators (\texttt{crbegin}/\texttt{crend} since C++11) & $O(1)$ & C++98 \\
  \addlinespace[3pt]
  \multicolumn{4}{@{}l}{\textbf{Capacity}}\\
  \addlinespace[2pt]
  \texttt{empty}\,$\dagger$ & empty test ($\texttt{begin()==end()}$); \texttt{[[nodiscard]]} since C++20 & $O(1)$ & C++98 \\
  \texttt{size} & number of elements & $O(1)$ & C++98 \\
  \texttt{max\_size} & theoretical maximum size & $O(1)$ & C++98 \\
  \addlinespace[3pt]
  \multicolumn{4}{@{}l}{\textbf{Modifiers}}\\
  \addlinespace[2pt]
  \texttt{clear}\,$\dagger$ & remove all elements & $O(n)$ & C++98 \\
  \texttt{insert}\,$\dagger$ & element(s) without overwriting (value, \texttt{P\&\&}, hint, range, \texttt{initializer\_list}, node handle since C++17) & $O(\log n)$; hint amort.\ $O(1)$; range $O(N\log(n{+}N))$ & C++98 \\
  \texttt{insert\_range} & insert an input range, no overwrite & $O(N\log(n{+}N))$ & C++23 \\
  \texttt{insert\_\allowbreak{}or\_\allowbreak{}assign}\,$\dagger$ & insert, or assign to an existing value; \texttt{bool} reports insertion & $O(\log n)$ & C++17 \\
  \texttt{emplace}\,$\dagger$ & construct in place, only if the key is absent (otherwise discarded) & $O(\log n)$ & C++11 \\
  \texttt{emplace\_hint} & like \texttt{emplace}, near the hint & $O(\log n)$; amort.\ $O(1)$ & C++11 \\
  \texttt{try\_emplace}\,$\dagger$ & like \texttt{emplace}, leaves the arguments untouched on an existing key & $O(\log n)$ & C++17 \\
  \texttt{erase}\,$\dagger$ & by iterator/range/key (key: 0 or 1 removed) & iter.\ amort.\ $O(1)$; key $O(\log n)$ & C++98 \\
  \texttt{swap} & exchange contents and comparator; iterators stay valid & $O(1)$ & C++98 \\
  \texttt{extract} & unlink an element as a node handle (no copy/move) & iter.\ amort.\ $O(1)$; key $O(\log n)$ & C++17 \\
  \texttt{merge} & splice nodes from a source map; existing keys remain there & $O(N\log(n{+}N))$ & C++17 \\
  \addlinespace[3pt]
  \multicolumn{4}{@{}l}{\textbf{Lookup}}\\
  \addlinespace[2pt]
  \texttt{find}\,$\dagger$ & iterator to the key, or \texttt{end()} & $O(\log n)$ & C++98 \\
  \texttt{count}\,$\dagger$ & number (0 or 1) & $O(\log n)$ & C++98 \\
  \texttt{contains}\,$\dagger$ & membership (\texttt{bool}) & $O(\log n)$ & C++20 \\
  \texttt{lower\_bound}\,$\dagger$ & first key $\ge$ \texttt{key} & $O(\log n)$ & C++98 \\
  \texttt{upper\_bound}\,$\dagger$ & first key $>$ \texttt{key} & $O(\log n)$ & C++98 \\
  \texttt{equal\_range}\,$\dagger$ & pair \texttt{[lower\_bound,\allowbreak\,upper\_bound)} & $O(\log n)$ & C++98 \\ % chktex 9
  \addlinespace[3pt]
  \multicolumn{4}{@{}l}{\textbf{Observers}}\\
  \addlinespace[2pt]
  \texttt{key\_comp}\,$\dagger$ & copy of the key comparator & $O(1)$ & C++98 \\
  \texttt{value\_comp} & \texttt{value\_compare} (compares pairs by key) & $O(1)$ & C++98 \\
  \addlinespace[3pt]
  \multicolumn{4}{@{}l}{\textbf{Non-member functions}}\\
  \addlinespace[2pt]
  \texttt{operator==} & equal size and pairwise-equal elements & $O(n)$ & C++98 \\
  \texttt{operator<=>} & lexicographical three-way comparison; synthesizes \texttt{<\,<=\,>\,>=} & $O(n)$ & C++20 \\
  \texttt{std::swap} & specialization; calls \texttt{lhs.swap(rhs)} & $O(1)$ & C++98 \\
  \texttt{std::erase\_if} & removes all elements with \texttt{pred==true}; returns the count & $O(n)$ & C++20 \\
  Deduction guides & CTAD from range/\allowbreak\texttt{initializer\_list}/\allowbreak\texttt{from\_range} (\texttt{from\_range} since C++23) & --- (compile time) & C++17 \\
  \bottomrule
\end{longtable}

\noindent The comparison operators \texttt{<\,<=\,>\,>=}, standalone before C++20, are synthesized
from \texttt{operator<=>} since C++20, and \texttt{!=} is rewritten from \texttt{operator==}. The
heterogeneous overloads of \texttt{operator[]}, \texttt{at}, \texttt{try\_emplace}, and
\texttt{insert\_or\_assign} (transparent key\footnote{Heterogeneous lookup since C++14: if the
comparator carries the type \texttt{is\_transparent}, the lookup functions accept arguments of a
type other than \texttt{key\_type} without forming a temporary key
object~\cite{wakely2013heterogeneous}; C++23 extends this to
\texttt{erase}/\texttt{extract}~\cite{boyarinov2021erasure}, C++26 to \texttt{operator[]},
\texttt{at}, \texttt{try\_emplace}, and
\texttt{insert\_or\_assign}~\cite{boyarinov2023heterogeneous}.}) belong only to C++26 and remain
excluded here; the C++23 additions are the \texttt{from\_range} constructor,
\texttt{insert\_range}, the transparent-key overloads of
\texttt{erase}/\texttt{extract}, and the associated \texttt{from\_range} deduction
guides~\cite{jabot2022ranges,boyarinov2021erasure}. The \texttt{constexpr} capability of
\texttt{std::map} likewise belongs only to C++26 and remains excluded~\cite{cppreference_map}.

\section{\texttt{std::vector} (Sequence)}\label{asec:if-vector}
Table~\ref{tab:if-vector} lists the complete interface of \texttt{std::vector<T,\,Allocator>} as of
C++23 ($n=\texttt{size()}$, for range operations $N$ the number of input elements; version markers
and complexities follow the standard text~\cite{iso_cpp} and are cross-checked against the online
reference~\cite{cppreference_vector}). All members are \texttt{constexpr} since
C++20\footnote{Since C++20, all member functions of \texttt{std::vector} are \texttt{constexpr}; a
\texttt{constexpr} object itself usually fails on the dynamic
storage~\cite{dionne2019constexprvector}.}. A $\dagger$ marks the core operations that the
conformance gate of the Sequence family checks before each measurement against a
\texttt{std::deque} oracle (append, index access, size, clear; for these four operations the
\texttt{deque} semantics coincide with those of \texttt{std::vector}): the sequence hull of this
work offers exactly these four functions across the module boundary, and index access is checked
there rather than undefined (status return: hit flag and output value at the module boundary,
\texttt{std::optional} in the organ). All other functions of the table --- capacity control,
position and range operations, \texttt{front}/\texttt{back}/\texttt{pop\_back}, \texttt{data},
iterators, \texttt{swap}, and the non-members --- do not cross the module boundary; growth at the
front end (\texttt{push\_front}), which the \texttt{deque} oracle would additionally offer, is
envisaged as an extension of the module boundary but not implemented, and is logged in the gate as
skipped, not as passed.

\begin{longtable}{@{}>{\raggedright\arraybackslash}p{3.1cm} >{\raggedright\arraybackslash}p{5.7cm} >{\raggedright\arraybackslash}p{2.7cm} >{\raggedright\arraybackslash}p{1.1cm}@{}}
  \caption{Complete \texttt{std::vector} interface (C++23)}\label{tab:if-vector}\\
  \toprule
  \textbf{Function} & \textbf{Official semantics} & \textbf{Complexity} & \textbf{Since} \\
  \midrule
\endfirsthead%
  \toprule
  \textbf{Function} & \textbf{Official semantics} & \textbf{Complexity} & \textbf{Since} \\
  \midrule
\endhead%
  \multicolumn{4}{@{}l}{\textbf{Construction, assignment, allocator}}\\
  \addlinespace[2pt]
  \texttt{(constructor)} & empty/\texttt{Allocator}; $n\times$ default/value; range; copy/move (incl.\ allocator-extended); \texttt{initializer\_list}; \texttt{from\_range} since C++23 & $O(1)$ or $O(n)$ & C++98 \\
  \texttt{(destructor)} & destroys all elements, frees storage & $O(n)$ & C++98 \\
  \texttt{operator=} & copy/move/\texttt{initializer\_list} & $O(n)$ & C++98 \\
  \texttt{assign} & replace contents with $n\times$ value/\allowbreak{}range/\allowbreak{}\texttt{initializer\_\allowbreak{}list} & $O(n)$ & C++98 \\
  \texttt{assign\_range} & replace contents with an input range & $O(n)$ & C++23 \\
  \texttt{get\_allocator} & copy of the associated allocator & $O(1)$ & C++98 \\
  \addlinespace[3pt]
  \multicolumn{4}{@{}l}{\textbf{Element access}}\\
  \addlinespace[2pt]
  \texttt{at}\,$\dagger$ & bounds-checked access; throws \texttt{std::out\_of\_range} & $O(1)$ & C++98 \\
  \texttt{operator[]} & unchecked access (UB out of range) & $O(1)$ & C++98 \\
  \texttt{front}/\texttt{back} & reference to first/last element (UB if empty) & $O(1)$ & C++98 \\
  \texttt{data} & pointer to the contiguous storage (not in \texttt{vector<bool>}) & $O(1)$ & C++11 \\
  \addlinespace[3pt]
  \multicolumn{4}{@{}l}{\textbf{Iterators}}\\
  \addlinespace[2pt]
  \texttt{begin}/\texttt{end}, \texttt{rbegin}/\texttt{rend} & forward/reverse iterators in index order (\texttt{c…} variants since C++11) & $O(1)$ & C++98 \\
  \addlinespace[3pt]
  \multicolumn{4}{@{}l}{\textbf{Capacity}}\\
  \addlinespace[2pt]
  \texttt{empty}/\allowbreak\texttt{size}\,$\dagger$/\allowbreak\texttt{max\_size} & empty test, element count, maximum size & $O(1)$ & C++98 \\
  \texttt{reserve} & capacity $\ge n$; reallocates and invalidates; throws \texttt{length\_error} if $n>$\texttt{max\_size()} & $O(n)$ & C++98 \\
  \texttt{capacity} & currently allocated capacity & $O(1)$ & C++98 \\
  \texttt{shrink\_to\_fit} & non-binding reduction of capacity toward \texttt{size()} & $O(n)$ & C++11 \\
  \addlinespace[3pt]
  \multicolumn{4}{@{}l}{\textbf{Modifiers}}\\
  \addlinespace[2pt]
  \texttt{clear}\,$\dagger$ & remove all elements (capacity unchanged) & $O(n)$ & C++98 \\
  \texttt{insert} & before a position: value/$n\times$/range/\allowbreak\texttt{initializer\_list}; shifts the tail & $O(n)$; range forms $O(n{+}N)$ & C++98 \\
  \texttt{insert\_range} & insert an input range before a position & $O(n{+}N)$ & C++23 \\
  \texttt{emplace} & construct an element in place before a position & $O(n)$ & C++11 \\
  \texttt{erase} & remove a position/range, shift the tail down & $O(n)$ & C++98 \\
  \texttt{push\_back}\,$\dagger$ & append an element (copy/move) & amort.\ $O(1)$ & C++98 \\
  \texttt{emplace\_back} & append in place; returns a reference since C++17 & amort.\ $O(1)$ & C++11 \\
  \texttt{append\_range} & append an input range & amort.\ $O(N)$ & C++23 \\
  \texttt{pop\_back} & remove the last element (UB if empty) & $O(1)$ & C++98 \\
  \texttt{resize} & truncate to $n$, or fill with default/value & $O(n)$ & C++98 \\
  \texttt{swap} & exchange contents and capacity with another vector & $O(1)$ & C++98 \\
  \addlinespace[3pt]
  \multicolumn{4}{@{}l}{\textbf{Non-member functions}}\\
  \addlinespace[2pt]
  \texttt{operator==} & equal size and element-wise equality & $O(n)$ & C++98 \\
  \texttt{operator<=>} & lexicographical three-way comparison & $O(n)$ & C++20 \\
  \texttt{std::swap} & specialization; calls \texttt{lhs.swap(rhs)} & $O(1)$ & C++98 \\
  \texttt{std::erase} & removes all elements equal to a value; returns the count & $O(n)$ & C++20 \\
  \texttt{std::erase\_if} & removes all elements with \texttt{pred==true}; returns the count & $O(n)$ & C++20 \\
  Deduction guides & CTAD from range/\allowbreak\texttt{from\_range} (\texttt{from\_range} since C++23) & --- (compile time) & C++17 \\
  \bottomrule
\end{longtable}

\noindent The specialization \texttt{std::vector<bool>} packs boolean values bitwise: it has no
\texttt{data()} and no contiguity guarantee, returns via \texttt{operator[]}/\allowbreak%
\texttt{at}/\allowbreak\texttt{front}/\allowbreak\texttt{back} a proxy
type\footnote{Stand-in object: \texttt{vector<bool>::reference} behaves in an expression like
\texttt{bool\&}, but denotes a single bit instead of an addressable object, which is why
\texttt{data()} and the contiguity guarantee are dropped~\cite{iso_cpp}.} \texttt{reference}
instead of \texttt{bool\&}, and additionally offers \texttt{flip()} (at the vector and at the
\texttt{reference} level), a static \texttt{swap(reference,\allowbreak\ reference)}, and a
\texttt{std::hash} specialization. The comparison operators behave as for \texttt{std::map}.

\section{Composite Functions and the Conformance Subset}\label{asec:if-decomp}
Several interface functions are internally compositions of more primitive operations of the same
interface (Table~\ref{tab:if-decomp}); the table describes the ISO standard contract. The
measurement hull of this work deliberately deviates from it in two places, because the ABI
conformance contract demands \texttt{noexcept} and a throwing hull counts as non-conformant:
\texttt{at} delivers a status return instead of the exception (hit flag and output value at the
module boundary, \texttt{std::optional} in the organ), and \texttt{operator[]} is not a call of its
own at the module boundary but is composed host-side from exact lookup and default insertion ---
with the ISO semantics (insertion on absence, no \texttt{out\_of\_range}), but without a reference
return across the binary boundary; the write path runs via \texttt{insert\_or\_assign}. In the
sequence hull, \texttt{operator[]} is moreover checked rather than undefined
(Section~\ref{asec:if-vector}). For this thesis that is the central observation: all map accessors
and modifiers of the hull rest on a single query primitive of the module boundary --- the exact
lookup, flanked by insertion and deletion ---; an axis implementation therefore only needs to
provide a correct exact lookup together with node linking, and the rest follows as a fixed
composition. The ordered operations --- \texttt{begin}/\texttt{end}, \texttt{lower\_bound},
\texttt{upper\_bound}, \texttt{equal\_range} --- additionally require an ordered search and are
therefore carried as a separate, additive capability of the module boundary; unordered axes
satisfy the same hull without it and pay for the missing order with linear effort on ordered
reading. Exactly this interchangeability is the object of the measurement
(Section~\ref{sec:measurement-system}). The conformance gate checks, per family, the core subset
marked with $\dagger$ in Tables~\ref{tab:if-map} and~\ref{tab:if-vector}.

\begin{longtable}{@{}>{\raggedright\arraybackslash}p{3.3cm} >{\raggedright\arraybackslash}p{5.6cm} >{\raggedright\arraybackslash}p{3.1cm}@{}}
  \caption{Composite functions (internal composition of primitive operations)}\label{tab:if-decomp}\\
  \toprule
  \textbf{Function} & \textbf{Composition} & \textbf{Note} \\
  \midrule
\endfirsthead%
  \toprule
  \textbf{Function} & \textbf{Composition} & \textbf{Note} \\
  \midrule
\endhead%
  \texttt{map::}\allowbreak\texttt{operator[]} & \texttt{lower\_bound} + default insertion on absence $\to$ reference & mutates on absence; no \texttt{out\_of\_range}, the insertion can throw \\
  \texttt{map::at} & \texttt{find} + throw \texttt{std::out\_of\_range} on absence & the only access that throws on absence; never mutates \\
  \texttt{map::insert} / \texttt{emplace} & ordered search + link a node only if the key is absent & no overwrite; \texttt{emplace} may still consume arguments on a hit \\
  \texttt{map::try\_\allowbreak{}emplace} & like \texttt{emplace}, but arguments untouched on a hit & preserves move-only arguments \\
  \texttt{map::insert\_\allowbreak{}or\_\allowbreak{}assign} & ordered search + insert or assign & overwrites (unlike \texttt{insert}) \\
  \texttt{map::count} / \texttt{contains} / \texttt{equal\_range} & reduce to \texttt{find} resp.\ \texttt{lower\_bound}+\texttt{upper\_bound} & always 0 or 1 for a unique-key map \\
  \texttt{vector::}\allowbreak\texttt{push\_back} / \texttt{emplace\_back} & possible growth (realloc + relocation) + construction at the end & hidden cost + invalidation on realloc \\
  \texttt{vector::insert} / \texttt{emplace} & possible growth + tail shift + insertion & extra $O(n)$ shift \\
  \texttt{vector::resize} & append (default/value) or truncate the tail & bidirectional; does not shrink capacity \\
  \texttt{vector::assign} & \texttt{clear} + insert range/$n\times$ & capacity reused where possible \\
  \bottomrule
\end{longtable}

\noindent At the module boundary, the order of primitiveness of this table is reversed: the one
insert drive of the SearchAlgorithm family (concept Map) carries the semantics of
\texttt{insert\_or\_assign} --- insert or
update, the return reports whether a new key arose ---, so that \texttt{insert\_or\_assign} is the
pass-through and the non-overwriting forms \texttt{insert}, \texttt{emplace}, and
\texttt{try\_emplace} are composed host-side as an exact lookup followed by insertion only on a
missing key: one path for three standard functions, which the gate sets against three oracle calls.
The ISO contract carries it the other way round (\texttt{insert} as the basic form,
\texttt{insert\_or\_assign} as the composition); both views are correct and describe different
levels --- exactly the contrast that makes the interchangeability behind a fixed interface
concrete.

\noindent Functions that both hulls \emph{share} by name --- \texttt{size}, \texttt{empty},
\texttt{clear}, the iterator families, \texttt{swap}, \texttt{get\_allocator}, and
\texttt{insert}/\texttt{emplace}/\texttt{erase} --- have the same role but different referents: key
semantics in the map, position/index semantics in the vector; \texttt{operator[]}, for instance,
inserts in the map, but is an unchecked index access in the vector --- checked in the sequence
hull of this work (Section~\ref{asec:if-vector}), composed host-side in the map hull of this work
(see above). This separation is structurally enforced in the cache engine: the Map and Sequence
families have separate drive interfaces and separate oracles, and an adapter foreign to its family is
ruled out by the type system. Which of these functions the module boundary provides, which the host
composes, and which are not wired up is recorded, per function with its hop path, by the developer
document named at the outset.
